\section{Les réunions internes}
\label{les-ruxe9unions-internes}

\subsection{L'objectif}
\label{lobjectif}

\begin{enumerate}
 \item
  La FCÉG organise quatre (4) réunions internes par an pour ses officiers, dans le but de coordonner les initiatives, la stratégie et les objectifs de la Fédération, ainsi que le processus de transition. Ces réunions, avec leurs objectifs spécifiques et leur calendrier, sont les suivantes :
\end{enumerate}

\subsection{Les réunions en ligne}
\label{les-ruxe9unions-en-ligne}

\subsubsection{Participation}
\label{participation}

\begin{enumerate}
 \item
  Tous les membres du conseil d\textquotesingle administration, les commissaires et les responsables d\textquotesingle activités.
 \item
  Les membres du groupe de travail sont invités à y participer.
\end{enumerate}

\subsubsection{La Réunion d\textquotesingle été:}
\label{la-ruxe9union-duxe9tuxe9}

\begin{enumerate}
 \item
  Accent mis sur les initiatives du groupe, avec un examen des progrès du plan d\textquotesingle action/stratégique.
 \item
  Elle sert de session de planification pour le sommet du développement des associations en ingénierie.
 \item
  Elle se tient traditionnellement à la fin du mois de juin ou en juillet et se déroule du vendredi au dimanche.
\end{enumerate}

\subsubsection{La Réunion d\textquotesingle automne:}
\label{la-ruxe9union-dautomne}

\begin{enumerate}
 \item
  Poursuite des travaux sur les initiatives des groupes, avec un examen plus approfondi des progrès du plan d\textquotesingle action stratégique.
 \item
  Elle sert de session de planification pour le CCLI
 \item
  Elle se tient traditionnellement à la fin du mois d\textquotesingle octobre ou au début du mois de novembre et se déroule du vendredi au dimanche.
\end{enumerate}

\subsection{Les réunions en personne/hybrides}
\label{les-ruxe9unions-en-personnehybrides}

\subsubsection{Participation}
\label{participation-1}

\begin{enumerate}
 \item
  Seuls les membres du conseil d\textquotesingle administration peuvent assister aux réunions, à l\textquotesingle exception des cas suivants:

  \begin{enumerate}
   \item
    Les Commissaires aux relations d\textquotesingle entreprise entrant.e et sortant.e peuvent participer à la réunion de printemps.
   \item
    Les responsables d\textquotesingle activité peuvent participer à la réunion de printemps conformément à l\textquotesingle accord d\textquotesingle activité.
  \end{enumerate}
 \item
  Des exceptions peuvent être accordées par le conseil d\textquotesingle administration pour permettre aux commissaires ou aux responsables d\textquotesingle activités de participer à la réunion de printemps.

  \begin{enumerate}
   \item
    Le/La vice-président.e responsable du commissaire ou du responsable d\textquotesingle activité doit fournir une lettre au conseil d\textquotesingle administration détaillant les tâches que l\textquotesingle individu accomplira lors de la réunion en personne et la raison pour laquelle il doit y participer.
   \item
    Ces lettres seront discutées et votées lors de la réunion du conseil d\textquotesingle administration qui suit la présentation de la lettre et nécessitent un vote à la majorité pour approuver la participation de l\textquotesingle individu.
  \end{enumerate}
 \item
  Ces réunions seront accessibles en ligne pour les membres du conseil d\textquotesingle administration qui ne sont pas en mesure d\textquotesingle y assister en personne et pour tout commissaire/responsable d\textquotesingle activité qui souhaite y participer.
\end{enumerate}

\subsubsection{La Réunion de printemps}
\label{la-ruxe9union-de-printemps}

\paragraph{Objectif}
\begin{enumerate}
 \item
  La réunion de printemps sert principalement de réunion de transition entre le conseil d\textquotesingle administration sortant et le conseil entrant de la Fédération.
 \item
  Du temps est accordé pour faciliter des réunions de transition spécifiques aux rôles, aux sous-groupes et au groupe dans son ensemble.
 \item
  Les initiatives de l\textquotesingle année précédente sont transmises à la nouvelle équipe et le processus de planification des actions est lancé.
 \item
  Le plan stratégique est examiné et, en fonction des progrès réalisés par l\textquotesingle équipe sortante, l\textquotesingle équipe entrante fixe ses objectifs pour l\textquotesingle année.
 \item
  La réunion de printemps sert également d\textquotesingle introduction à Ingénieurs Canada pour l\textquotesingle équipe entrante, une journée de la réunion de printemps étant organisée par Ingénieurs Canada dans leurs bureaux à Ottawa.
 \item
  La journée Ingénieurs Canada est planifiée par Ingénieurs Canada et comprend généralement~:

  \begin{enumerate}
   \item
    Une introduction à Ingénieurs Canada~;
   \item
    Un aperçu de leurs priorités pour l\textquotesingle année~;
   \item
    Une mise à jour sur les progrès de la FCEG par le conseil sortant~;
   \item
    Un aperçu des objectifs du conseil entrant~; et
   \item
    D\textquotesingle autres séances portant sur les activités d\textquotesingle Ingénieurs Canada, telles que le BCAP et les Voies d\textquotesingle accès au génie.
  \end{enumerate}

\end{enumerate}

\paragraph{Planification}
\begin{enumerate}
 \item
  La planification de la réunion de printemps relève de la responsabilité des président.e.s entrant.e et sortant.e ainsi que des VPFA entrant.e et sortant.e.
 \item
  La planification de la réunion de printemps doit commencer au plus tard 1 mois après le CCLI.
 \item
  La réunion de printemps ne peut être organisée avant la troisième semaine de mai afin de permettre l\textquotesingle utilisation de résidences universitaires peu coûteuses comme hébergement.
 \item
  Les dates de la réunion de printemps doivent être arrêtées au plus tard la première semaine de mars et doivent être basées sur les disponibilités des équipes entrante et sortante, l\textquotesingle équipe entrante ayant la priorité.
 \item
  Les participant.e.s doivent contacter leur faculté pour un financement de voyage au plus tard 1 mois après le CCLI.

  \begin{enumerate}
   \item
    Le non-respect de cette obligation entraînera leur inéligibilité au financement de voyage.
  \end{enumerate}

 \item
  Les hébergements doivent être réservés avant la deuxième semaine de mars.
 \item
  Le calendrier préliminaire doit être envoyé à tou.te.s les participant.e.s avant le 1er avril et tou.te.s les membres des conseils entrant et sortant doivent avoir la possibilité de donner leur avis sur le calendrier.
 \item
  Le calendrier définitif doit être envoyé à tou.te.s les participant.e.s avant le 1er mai.
 \item
  Tou.te.s les participant.e.s qui en ont besoin doivent soumettre leur demande de financement de voyage FCEG avant le 1er mars.
 \item
  Les estimations initiales lors du processus de planification doivent prévoir 20 participant.e.s.
\end{enumerate}

\subsubsection{La pré-semaine du CCLI:}
\label{la-pruxe9-semaine-du-ccli}

\begin{enumerate}
 \item
  Mener à terme les initiatives du groupe et le plan d\textquotesingle action pour les progrès réalisés par rapport au plan stratégique.
 \item
  Se familiariser avec le lieu et le CO de l'activité
 \item
  Finaliser la liste des représentants des membres du CCLI et commencer à planifier la transition.
 \item
  Cette semaine n'a pas comme objectif pour les membres de l'exécutif national, y compris le/la VPS, d'exercer les fonctions des responsables d'activité et du CO de l'activité tel que défini dans leur accord d'activité respectif.
 \item
  Elle se déroule selon le calendrier suivant, en tenant compte des besoins concurrents de la période des fêtes et de l\textquotesingle agenda chargé du CCLI:

  \begin{enumerate}
   \item
    si le CCLI commence le 2 janvier, elle se tient du 29 décembre au 2 janvier;
   \item
    si le CCLI commence le 3 janvier, elle se tient du 30 décembre au 3 janvier;
   \item
    et si le CCLI commence le 4 janvier ou plus tard, la pré-semaine du CCLI doit commencer trois jours avant le CCLI lui-même:
  \end{enumerate}

 \item
  Il appartient au Conseil d\textquotesingle administration, après consultation du/de la VPS de la FCÉG et du/de la responsable d'activité du CCLI, de réduire ou d\textquotesingle augmenter la durée de la pré-semaine en fonction du nombre d\textquotesingle éléments en suspens à finaliser.
 \item
  Le responsable d\textquotesingle activité du CCLI est informé de la décision finale concernant la durée de la pré-semaine au plus tard une semaine après le SDAI tenue avant le CCLI.
\end{enumerate}

\subsubsection{Couverture des dépenses}
\label{couverture-des-depenses}

\paragraph{Meilleures pratiques}
\begin{enumerate}
 \item
  Les individus doivent agir de bonne foi et être conscients de l\textquotesingle utilisation des cotisations des membres.
 \item
  Les individus ne doivent pas chercher à dépenser le maximum de l\textquotesingle allocation possible qui leur est accordée.
 \item
  Tout achat jugé excessif par le/la VPFA devra faire l\textquotesingle objet d\textquotesingle une discussion entre l\textquotesingle individu, le/la président.e, les conseiller.ères nationaux.ales et le/la VPFA.
 \item
  Cela peut également avoir un impact sur le remboursement de l\textquotesingle individu.
\end{enumerate}

\paragraph{Allocation de repas}
\begin{enumerate}
 \item
  La FCEG couvre jusqu\textquotesingle à un montant défini par repas pour les personnes assistant aux réunions internes en personne.
 \item
  Pour recevoir l\textquotesingle allocation pour un repas donné, le membre de l\textquotesingle équipe doit être présent à la réunion pour ce repas.
 \item
  Les allocations de repas sont des prix avant taxes et avant pourboire.
 \item
  La FCEG remboursera les individus à la fin de la réunion en fonction de leurs dépenses réelles.
 \item
  Les individus qui dépensent moins que l\textquotesingle allocation totale qui leur est accordée ne seront remboursés que du montant qu\textquotesingle ils ont dépensé.
 \item
  Les individus qui dépensent plus que l\textquotesingle allocation totale qui leur est accordée ne seront remboursés que du montant total de l\textquotesingle allocation.
 \item
  Les boissons gazeuses, les jus, l\textquotesingle eau et les boissons alcoolisées ne peuvent pas être couverts par l\textquotesingle allocation de repas.
 \item
  Les membres sont encouragés à apporter leurs propres bouteilles d\textquotesingle eau réutilisables à toutes les réunions internes.
 \item
  Le dessert ne peut pas être couvert par l\textquotesingle allocation de repas.
 \item
  Si un repas est fourni par un partenaire, aucune allocation de repas ne sera accordée pour ce repas.
 \item
  L\textquotesingle allocation de repas pour chaque repas est définie ci-dessous~:

  \begin{enumerate}
   \item
    10~\$ Petit-déjeuner
   \item
    10~\$ Déjeuner
   \item
    15~\$ Dîner
  \end{enumerate}

\end{enumerate}

\paragraph{Transport lors des réunions internes}
\begin{enumerate}
 \item
  Le transport ne peut être couvert que si la distance entre les hébergements et les lieux de réunion est jugée déraisonnable à pied par le/la VPFA et le/la président.e.
 \item
  Des exceptions au cas par cas peuvent être accordées à la discrétion du/de la VPFA et du/de la président.e.
 \item
  Lorsque le transport est jugé justifié, les individus ne peuvent être remboursés que pour l\textquotesingle option de transport la moins chère et réalisable, à la discrétion du/de la VPFA et du/de la président.e.
 \item
  Les individus qui choisissent d\textquotesingle utiliser une option de transport plus coûteuse ne seront pas remboursés pour l\textquotesingle ensemble du coût.
\end{enumerate}

\subsection{La logistique}
\label{la-logistique}

\begin{enumerate}
 \item
  Ces réunions sont organisées par l'exécutif national.
 \item
  Le budget principal de la FCÉG sera alloué pour couvrir les frais d\textquotesingle hébergement, de restauration et de logistique des réunions de printemps, d\textquotesingle été et d\textquotesingle automne de la FCÉG.
 \item
  Il sera également alloué pour couvrir la nourriture et la logistique de la pré-semaine du CCLI, tandis que l\textquotesingle hébergement sera couvert par le budget du/de la responsable d\textquotesingle activité du CCLI.
 \item
  En outre, l\textquotesingle exécutif national organise une réunion de transition pour l\textquotesingle exécutif national et les personnes qui ont été élues à l\textquotesingle exécutif national pour l\textquotesingle année suivante.
 \item
  Cette réunion est organisée par l\textquotesingle exécutif national en place et a lieu après le CCLI, où se déroulent les élections, et avant la réunion de printemps.
\end{enumerate}

